\section{Infinite-mode truncation and error propagation}
\label{tail:section}\label{sec:tails}

This section supplies analytic estimates for the modes omitted after
$m=38$.  We use the radial functions $b_m,U_m,u_m,B,D_m$ of the preceding
construction, with
\[
 \begin{gathered}
 a=\sqrt{2/\sqrt3},\qquad p=a/2,\qquad R_0=a/\sqrt3,\\
 b=\frac{28209}{100000},\qquad h_s=\frac{3}{40},\qquad
 c=\frac{3}{50},\qquad l=R_0+\frac{19}{250}.
 \end{gathered}
\]
Thus $k_0=6$, $k_m=6m$ for $m\geq1$, and
\[
 c_m=\frac1{k_m}\sum_{z\in\Lambda_\triangle\setminus\{0\}}z^{-k_m}.
\]
Tildes will always mean that the modes $m\geq39$ have been omitted;
the index-zero term is retained.  Differences between a full quantity
and its truncation are denoted by $\Delta$.
Our goal is an absolute error below $10^{-11}$ for $f_0,e$ and their
first two derivatives on $[b,l]$, conditional on the retained-mode
bounds verified in Section~\ref{sec:computation}.

\subsection{Coefficient estimates}
\label{tail:coefficients}

Write a lattice point as
$z=a(n+q/2+i\sqrt3 q/2)$, with $(n,q)\in\Z^2$.  Then
\[
 |z/a|^2=n^2+nq+q^2\geq\frac34\max(|n|,|q|)^2.
\]
There are $8v$ integer pairs with $\max(|n|,|q|)=v$.  For every
$d=6m\geq6$ and every positive integer $D$, the sum outside the square
$|n|,|q|\leq D$, before multiplication by $a^{-d}$, is therefore at most
\begin{equation}
 8(4/3)^{d/2}\sum_{v>D}v^{1-d}
 \leq 8(4/3)^{d/2}\frac{D^{2-d}}{d-2}.
 \label{tail:coefficient-square}
\end{equation}
The last inequality follows by comparison with the integral from $D$
to infinity.  The first shell has six points with squared norm $1$
and two with squared norm $3$.  Moreover,
\[
 \sum_{v\geq2}v^{1-d}
 \leq 2^{1-d}+\int_2^\infty x^{1-d}\,\dd x
 =2^{1-d}\left(1+\frac2{d-2}\right)
 \leq3\,2^{1-d}.
\]
It follows that, for $d=k_m\geq12$,
\begin{equation}
 |c_m|\leq\frac{1}{d a^d}
 \left(6+2\,3^{-d/2}+24(4/3)^{d/2}2^{1-d}\right)
 \leq\frac7{d a^d}.
 \label{tail:coefficient-uniform}
\end{equation}
Indeed, the bracket equals $6+50/729<7$ at $d=12$, and its two
nonconstant terms decrease with $d$.  The finite coefficient
enclosures may use \eqref{tail:coefficient-square} with $D=256$ for
$m=1$ and $D=48$ for $2\leq m\leq38$, including the factor
$1/(k_m a^{k_m})$.  Reflection symmetry makes each coefficient real.

\subsection{Uniform derivative bounds for omitted modes}
\label{tail:radial-derivatives}

For $k=6m\geq234$, define the rational number
\begin{equation}
 w_k=10^6\left(
 30k^2(13/20)^k+3\cdot10^6 k^3 2^{-k}
 +5600k^5(3/5)^k+7\cdot10^{11}k^2(289/500)^k
 \right).
 \label{tail:w-definition}
\end{equation}
For $0\leq j\leq3$ and $0<r\leq l$, the construction satisfies
\begin{equation}
 |u_m^{(j)}(r)|\leq w_k,
 \qquad |D_m^{(j)}(r)|\leq w_k.
 \label{tail:mode-derivative-bound}
\end{equation}
At spline knots these inequalities mean both one-sided derivatives.
Here are explicit checks of the constants.

The polynomial cutoff $S(y)=10y^3-15y^4+6y^5$ on $[0,1]$, extended
constantly outside that interval, obeys
\[
 |S|\leq1,\qquad |S'|\leq2,\qquad
 |S''|\leq6,\qquad |S'''|\leq60.
\]
The identities $p^4=1/12$ and $R_0^4=4/27$ give the convenient
rational bounds
\[
 .537<p<.538,\quad .620<R_0<.621,\quad .696<l<.697,
 \quad l/a<13/20,\quad R_0/a<289/500.
\]
All terminating decimals in this section denote rational numbers.

For
$U_m(r)=(-1)^m c_mr^k(1-S((r-R_0)/.1))$,
Leibniz' rule and \eqref{tail:coefficient-uniform} bound the third
derivative by
\[
 7k^2(l/a)^k
 \left(l^{-3}+\frac{60}{kl^2}
             +\frac{1800}{k^2l}+\frac{60000}{k^3}\right).
\]
Using $l>.69$ and $k\geq234$, the coefficient of
$k^2(13/20)^k$ is less than $25.445<30$.
The same calculation for orders $0,1,2$ gives coefficients below
$0.000001$, $0.000197$, and $0.070572$, respectively.  This proves the
first term inside the parentheses in \eqref{tail:w-definition}.

For $p-c\leq r\leq p$, write
$q_m(r)=T_{0m}+T_{1m}(r-p)+T_{2m}(r-p)^2$, so that
$b_m(r)=S((r-p+c)/c)q_m(r)$.  Since $p/a=1/2$ and $p>1/2$,
\[
 |T_{0m}|\leq\frac7k2^{-k},\qquad
 |T_{1m}|\leq14k2^{-k},\qquad
 |T_{2m}|\leq\frac{14}{3}k^3 2^{-k}.
\]
The last inequality uses
$((k-2)(k-3)/6)-1/2\leq k^2/6$.
Consequently
\begin{align*}
 |q_m|&\leq(7/k+.84k+.0168k^3)2^{-k},\\
 |q_m'|&\leq(14k+.56k^3)2^{-k},\\
 |q_m''|&\leq(28/3)k^3 2^{-k}.
\end{align*}
In particular,
\[
 |b_m'''|\leq\frac{60}{c^3}|q_m|
       +\frac{18}{c^2}|q_m'|+\frac6c|q_m''|
 <8500k^3 2^{-k}<3\cdot10^6 k^3 2^{-k}.
\]
The lower derivative orders satisfy the same final bound, and below
$p-c$ the function is zero.

On $[p,p+.01]$, put $w=r^2-p^2$ and $X=p+\sqrt w$.  The polynomial
formula for $V_m$ gives, for $0\leq\ell\leq3$,
\[
 \left|\frac{\dd^\ell V_m}{\dd w^\ell}\right|
 \leq |c_m| k^{2\ell}X^{k-2\ell}.
\]
For $j\geq\ell$, the coefficient comparison that proves this is
\[
 j(j-1)\cdots(j-\ell+1)\binom{k}{2j}
 \leq k^{2\ell}\binom{k-2\ell}{2j-2\ell}.
\]
The bound $p>.537$ implies
$.02p+.0001<.04p^2$, hence $X\leq1.2p=(3/5)a$ on this interval.
Also $r<.55$ and $X>1/2$.  The chain rule gives
\[
 V_m'''(r)=8r^3\frac{\dd^3V_m}{\dd w^3}
                     +12r\frac{\dd^2V_m}{\dd w^2}.
\]
Thus its absolute value is at most
\[
 7\left(8(.55)^3 2^6+
           \frac{12(.55)2^4}{234^2}\right)k^5(3/5)^k
 <597k^5(3/5)^k.
\]
The lower derivatives satisfy smaller bounds.  The common bound
$5600k^5(3/5)^k$ therefore applies through order three.

On $[p+.01,R_0]$, use $z=p+i\sqrt{r^2-p^2}$, so that
$V_m=\Re(c_mz^k)$ and $|z|=r$.  Direct differentiation gives
\[
 |z'|\leq7,\qquad |z''|\leq300,\qquad |z'''|\leq60000.
\]
For example, $\sqrt{r^2-p^2}>.103$, $p<.54$, and $r<.63$ suffice
in the formulas $|z'|=r/\sqrt{r^2-p^2}$,
$|z''|=p^2/(r^2-p^2)^{3/2}$, and
$|z'''|=3p^2r/(r^2-p^2)^{5/2}$.
Since $1/r<2$, the third derivative of $c_mz^k$ is bounded by
\[
 7\left(8\cdot343+\frac{12\cdot7\cdot300}{234}
                    +\frac{2\cdot60000}{234^2}\right)
 k^2(289/500)^k
 <20000 k^2(289/500)^k.
\]
All derivative orders through three are therefore bounded by
$3\cdot10^6 k^2(289/500)^k$ on this interval.

For the extension
$\mathcal HV_m=V_m(R_0)+P(r-R_0)V_m'(R_0)+Q(r-R_0)V_m''(R_0)$,
direct bounds on the coefficients of $P,Q$ show, for $0\leq j\leq3$,
\[
 |P^{(j)}|\leq2\cdot10^5,\qquad |Q^{(j)}|\leq4000.
\]
Thus the common bound for the extension is at most
\[
 (1+200000+4000)\,3\cdot10^6 k^2(289/500)^k
 <7\cdot10^{11} k^2(289/500)^k.
\]
It remains valid on the constant part beyond $R_0+h_s$.

Finally the cutoff
$\tau(r)=S((r-R_0-.001)/.075)$ satisfies
\[
 |\tau^{(j)}|\leq1,27,1067,142223
 \qquad(j=0,1,2,3).
\]
The maximal Leibniz factor through order three is
$1+3\cdot27+3\cdot1067+142223=145506<10^6$.
Applying this to
$u_m=(1-\tau)U_m+\tau b_m$ and
$D_m=(1-\tau)(b_m-U_m)$ proves
\eqref{tail:mode-derivative-bound}.  The series and their indicated
derivatives consequently converge absolutely and locally uniformly.
For $r<b$, the omitted $u_m$ derivatives additionally have their
original $O(r^{k-3})$ behavior.

\subsection{Norms and bounds for retained modes}
\label{tail:norms}

For a function $F$ with one-sided derivatives as needed, define
\[
 |F|_{j,t}=\sum_{q=0}^j\frac{|F^{(q)}(t)|}{q!},
 \qquad
 \|F\|_{j,I}=\sup_{t\in I}|F|_{j,t}.
\]
These norms are submultiplicative.  For nonnegative weights
$\omega_i$, a weighted mode norm means
$\sup_{t\in I}\sum_i\omega_i|F_i|_{j,t}$:
the sum is taken before the supremum.

The finite interval verification in Section~\ref{sec:computation}
establishes the following bounds for the retained modes:
\begin{equation}
 |\widetilde B_1|_{2,t}<6000,
 \qquad
 \sum_{i=0}^{38}(k_i+1)^3|\widetilde D_i|_{3,t}<2\cdot10^6,
 \qquad b\leq t\leq l.
 \label{tail:finite-core-premise}
\end{equation}
The finite coefficient enclosure also supplies $|c_1|<1$.
These assertions are not conclusions of the analytic estimates in
this section.  Conditional on them, the remaining estimates below
are independent of the finite calculation.

On $[0,b]$,
\[
 \widetilde B_1(t)=\pi t+
       \sum_{m=1}^{38}(-1)^m k_mc_mt^{k_m-1}
\]
has norm $|\widetilde B_1|_{2,t}<10$.  For an explicit check, use
$t<b<2/7$, $a>1$, $\pi<22/7$, and $|c_1|<1$ for the first mode.
The modes $k\geq12$ contribute at most
$\sum_{k=12,18,\ldots}7k^2(2/7)^{k-3}<.014$;
the first term and its geometrically decreasing ratios prove this
last rational inequality.  Together with the linear term and the
$k=6$ term this is less than $10$.

\subsection{Summing the omitted modes}
\label{tail:sum}

Set
\begin{equation}
 \mathfrak a=\sum_{k=234,240,\ldots}(k+1)^4w_k
       <2\cdot10^{-21},\qquad
 W_*=4\sum_{k=234,240,\ldots}w_k\leq4\mathfrak a.
 \label{tail:weighted-sum}
\end{equation}
Here is a rational verification of the numerical inequality.  For
each term $(k+1)^4k^sr^k$, with
\[
 (s,r)=(2,13/20),(3,1/2),(5,3/5),(2,289/500),
\]
the ratio at the next allowed index is
\[
 \left(\frac{k+7}{k+1}\right)^4
 \left(\frac{k+6}{k}\right)^s r^6.
\]
It decreases with $k$, and at $k=234$ it is strictly below $1/10$
for all four pairs.  The first weighted summand obeys the exact
rational comparison
\[
 \begin{split}
 235^4\,10^6\bigl(&30\cdot234^2(13/20)^{234}
 +3\cdot10^6\cdot234^3 2^{-234}\\
 &+5600\cdot234^5(3/5)^{234}
 +7\cdot10^{11}\cdot234^2(289/500)^{234}\bigr)
 <1.56\cdot10^{-21}.
 \end{split}
\]
Multiplying by $10/9$ gives a bound less than
$2\cdot10^{-21}$, as asserted in \eqref{tail:weighted-sum}.

Since $\Delta B_1=\sum_{m\geq39}u_m'$, the derivative estimates give
\[
 |\Delta B_1|_{2,t}\leq W_*,\qquad 0\leq t\leq l,
\]
and, on $[b,l]$,
\[
 \sum_{i\geq39}(k_i+1)^j|D_i|_{3,t}\leq4\mathfrak a,
 \qquad 0\leq j\leq4.
\]
Indeed the divided-factorial sums use factors $5/2$ and $8/3$,
respectively, both less than $4$.

\subsection{Propagation to the sector functions}
\label{tail:propagation}

Recall the quantities entering the sector formula:
\begin{align*}
 B_1(t)&=B'(t)+1/t,&
 S_1(t)&=\int_0^t B(r)r\,\dd r,\\
 S_2(t)&=\int_0^t(-2B_1(r)+rB_1(r)^2)\,\dd r,&
 H_k(t)&=t^{-k}\int_0^t B_1(r)r^k\,\dd r,\\
 Q_i(t)&=tD_i'(t)-k_iD_i(t).
\end{align*}
Write $N=2\cdot10^6$ temporarily.
Direct differentiation of $\Delta S_1$ gives
\[
 |\Delta S_1|_{2,t}\leq3W_*.
\]
For example $|\Delta B^{(j)}|\leq W_*/4$ through order three,
so the norm is at most
$[t^2/2+t+(1+t)/2]W_*/4$ for $t\leq l<1$.
For the other integral, use
\[
 \Delta(-2B_1+tB_1^2)
 =\bigl(-2+t(2\widetilde B_1+\Delta B_1)\bigr)\Delta B_1.
\]
The second-order norm of its integral is at most twice the supremum
of the first-order norm of the integrand, because $l<1$.
Submultiplicativity and $|t|_{1,t}<2$ yield
\[
 |\Delta S_2|_{2,t}
 \leq(4+8\cdot6000+4W_*)W_*<60000W_*.
\]

The identities and elementary bounds
\[
 H_k'=B_1-kH_k/t,\quad
 |H_k|\leq\frac{t\|B_1\|_{0,[0,t]}}{k+1},\quad
 |H_k'|\leq2\|B_1\|_{0,[0,t]},\quad
 |H_k''|\leq|B_1'|+\frac{2k}{t}\|B_1\|_{0,[0,t]}
\]
imply, for $b\leq t\leq l$,
\begin{equation}
 |H_k|_{2,t}\leq6(k+1)(6000+W_*),\qquad
 |\Delta H_k|_{2,t}\leq6(k+1)W_*.
 \label{tail:H-estimate}
\end{equation}
Indeed $l/(k+1)+5/2+k/b<6(k+1)$.  The bound for $\Delta H_k$
follows by the same linear argument applied to $\Delta B_1$.

The radial source and energy functions are
\begin{align*}
 f_0(t)={}&-2\pi\left(S_1(t)+t^2\sum_i\frac{D_i(t)}{k_i+2}\right)\\
 &-\frac12\left(\log t+S_2(t)
 +\sum_i(-2D_i+2k_iD_iH_{k_i})(t)
 +\sum_{i,j}\frac{k_ik_j}{k_i+k_j}D_i(t)D_j(t)\right),\\
 e(t)={}&\frac12\sum_{i,j}\frac{Q_i(t)Q_j(t)}{k_i+k_j}.
\end{align*}
The preceding bounds give
\begin{align}
 |f_0-\widetilde f_0|_{2,t}
 &\leq2\pi(3W_*+16\mathfrak a)\notag\\
 &\quad+\frac12\bigl(60000W_*+8\mathfrak a
       +12(2\cdot10^6)W_*
       +48(6000+W_*)\mathfrak a\notag\\
 &\hspace{44mm}+8(2\cdot10^6)\mathfrak a
       +16\mathfrak a^2\bigr)
 <6\cdot10^7\mathfrak a.
 \label{tail:f-estimate}
\end{align}
For completeness, $|t^2|_{2,t}=(t+1)^2<4$ gives the $16\mathfrak a$ source
term.  The sum of $-2D_i$ contributes $8\mathfrak a$.  Split the mixed term as
\[
 \sum_i k_i(D_iH_{k_i}-\widetilde D_i\widetilde H_{k_i})
 =\sum_{i\leq38}k_i\widetilde D_i\Delta H_{k_i}
    +\sum_{i\geq39}k_iD_iH_{k_i}.
\]
The finite and omitted weighted norms give $12NW_*$ and
$48(6000+W_*)\mathfrak a$ after including its factor $2$.
Bounding $k_ik_j/(k_i+k_j)$ by $k_ik_j$, the quadratic difference
is at most $8N\mathfrak a+16\mathfrak a^2$.  These are precisely the terms in
\eqref{tail:f-estimate}.  Finally $W_*\leq4\mathfrak a$, $\pi<22/7$, and
$\mathfrak a<2\cdot10^{-21}$ make its total coefficient less than
$56,264,181<6\cdot10^7$.

Since $k_i\geq6$ and $t<1$, direct differentiation of $Q_i$ shows
\[
 |Q_i|_{2,t}\leq2(k_i+1)|D_i|_{3,t}.
\]
For example, its coefficients on
$|D_i|,|D_i'|,|D_i''|,|D_i'''|$ are bounded respectively by
$k_i$, $t+k_i-1$, $t+(k_i-2)/2$, and $t/2$.
The finite and omitted sums of these norms are bounded by $2N$
and $8\mathfrak a$.  Dropping the denominator $k_i+k_j\geq1$ in the energy
difference therefore gives
\begin{equation}
 |e-\widetilde e|_{2,t}
 \leq16(2\cdot10^6)\mathfrak a+32\mathfrak a^2.
 \label{tail:e-estimate}
\end{equation}
The norm contains half the absolute second derivative.  Accounting
for this factor, \eqref{tail:f-estimate} and
\eqref{tail:e-estimate} imply absolute truncation errors less than
$10^{-11}$ for every member of
\begin{equation}
 f_0,\quad e,\quad f_0',\quad e',\quad f_0'',\quad e''
 \qquad (b\leq t\leq l).
 \label{tail:uniform-error}
\end{equation}
In fact the respective second derivative bounds are below
$2.4\cdot10^{-13}$ and $1.3\cdot10^{-13}$.

\subsection{The constant region}
\label{tail:constant-region}

For $r\geq l$, all $D_i$ vanish and $B$ is constant.  The omitted
constant satisfies
\[
 |\Delta B(l)|\leq
 \sum_{k=234,240,\ldots}\frac7k(R_0/a)^k
 \leq\frac{(7/234)(289/500)^{234}}{1-(289/500)^6}
 <6.075\cdot10^{-58}<10^{-11}.
\]
Since $B_1(t)=1/t$ there, the derivatives of $\log t$ and $S_2(t)$
cancel.  The full infinite-mode functions consequently obey the
exact formula
\begin{equation}
 e(t)=0,\qquad
 f_0(t)=f_0(l)-\pi B(l)(t^2-l^2),\qquad t\geq l.
 \label{tail:exterior-formula}
\end{equation}
Thus outside $[b,l]$ one propagates the enclosures of $f_0(l)$ and
$B(l)$ through this exact formula; no uniform absolute error for
arbitrarily large $t$ is asserted.
